\chap{Review Laboratory: Circuits with Diodes and Transistors}

\section{Introduction}
This laboratory is a review of the first five. It contains typical problems of a mid-term project: equivalent resistance, Y--$\Delta$ transformation, circuits with diodes and Zener diodes, and circuits with bipolar transistors. You can use it to prepare for the mid-term examination.

\textbf{After this laboratory you will be able to:}
\begin{itemize}
  \item apply Y--$\Delta$ transformations and verify them by simulation;
  \item analyse circuits with diodes and Zener diodes by deciding the conduction state;
  \item analyse BJT circuits and detect saturation with a DC sweep;
  \item combine Zener diodes and transistors in a regulator and study its load regulation;
  \item use \cmd{.step}, \cmd{mc()} and \cmd{.tf} to automate the analysis.
\end{itemize}

\textbf{Exercises.} Save every schematic as \file{Lab6/ExNN.asc} (\file{Ex01.asc} to \file{Ex10.asc} for main exercises; \file{Ex11.asc} and \file{Ex12.asc} for advanced exercises). Calculation spaces are blank; fill in the comparison tables with your own theoretical and simulated values.

\section{LTspice guide}

\subsection{Strategy for a review laboratory}
This laboratory brings together the previous ones: resistor networks, diodes, Zener diodes and bipolar transistors. For each exercise, first decide the operating state of every non-linear device (is the diode conducting? is the Zener in breakdown? is the transistor active or saturated?) and write down your assumptions, then calculate, then simulate. A wrong assumption shows up as an inconsistent result (for example a negative diode current).

\subsection{Useful LTspice commands}
\textbf{DC sweep}: \cmd{.dc V1 0 6 0.05}. \textbf{Parameters and steps}: \cmd{.param RL=100}, \cmd{.step param RL list 50 100 1k}. \textbf{Transient}: \cmd{.tran 0 4m 0 5u}. \textbf{Measurements}: \cmd{.meas TRAN Vmax MAX V(out)}. \textbf{Random values}: \cmd{mc(value,tolerance)} used with \cmd{.step param run 1 20 1}. \textbf{Thevenin equivalent}: \cmd{.tf V(out) V1}.

\subsection{Checking a bias point}
After a \cmd{.op} run, check each non-linear device: a conducting diode has about 0.6--0.8\,V across it; a Zener in breakdown has a voltage close to $V_Z$ and a reverse current; a transistor in the active region has $V_{CE}>0.2$\,V and $I_C\approx\beta I_B$.

\section{Exercises}

\subsection{Exercise 1: Equivalent resistance of a bridge network}
\status{Modified (Ch.~7, Ex.~1: bridge network equivalent resistance via Y--$\Delta$ transformation).}
In the circuit of Figure~\ref{fig:lab6_ex01} a 24\,V source feeds a 4\,$\Omega$ resistor in series with a bridge made of 20, 30, 60 and 40\,$\Omega$ resistors and a 10\,$\Omega$ resistor between the middle nodes. Use a Y--$\Delta$ transformation to find the total equivalent resistance seen by the source, the current delivered by the source and the power dissipated in the 10\,$\Omega$ resistor. Then simulate.

\labfig{lab6_ex01.png}{10cm}{Exercise 1: bridge network}{fig:lab6_ex01}

\theory{10}

\simulation{
\item Draw the circuit of Figure~\ref{fig:lab6_ex01} in LTspice and save it as \file{Lab6/Ex01.asc}.
\item Add the \cmd{.op} analysis (\textit{Simulate $\rightarrow$ Edit Simulation Cmd $\rightarrow$ DC op pnt}), run the simulation and read the node voltages and branch currents.
\item $R_{eq}=V_1/I$, where $I$ is the magnitude of the source current.
\item Compute the power in the 10\,$\Omega$ resistor from its voltage and current.
\item Be careful with signs: LTspice reports the current of a voltage source as positive when it flows \emph{into} its positive terminal, and the current of a two-terminal device as positive when it flows from its first pin to its second pin.
}

\comparetable{
\cmprow{$R_{eq}$ ($\Omega$)}
\cmprow{$I$ (A)}
\cmprow{$P_{10\,\Omega}$ (W)}
}

\commentlines{2}

\subsection{Exercise 2: Equivalent resistance with equal resistors}
\status{Modified (Ch.~7, Ex.~2: balanced bridge network equivalent resistance).}
Five equal resistors R = 1\,k$\Omega$ are connected as in Figure~\ref{fig:lab6_ex02}: two in the upper branch, two in the lower branch and one in the middle between the two intermediate nodes. Find the equivalent resistance between A and B and the current in the middle resistor, and explain the result (what is the condition for a balanced bridge?). A 10\,V source is connected between A and B for the simulation.

\labfig{lab6_ex02.png}{8cm}{Exercise 2: network of equal resistors}{fig:lab6_ex02}

\theory{8}

\simulation{
\item Draw the circuit of Figure~\ref{fig:lab6_ex02} in LTspice and save it as \file{Lab6/Ex02.asc}.
\item Add the \cmd{.op} analysis (\textit{Simulate $\rightarrow$ Edit Simulation Cmd $\rightarrow$ DC op pnt}), run the simulation and read the node voltages and branch currents.
\item Read the current of the middle resistor and the voltage across it.
\item Be careful with signs: LTspice reports the current of a voltage source as positive when it flows \emph{into} its positive terminal, and the current of a two-terminal device as positive when it flows from its first pin to its second pin.
}

\comparetable{
\cmprow{$R_{AB}$ (k$\Omega$)}
\cmprow{Current in the middle resistor (mA)}
\cmprow{Voltage across the middle resistor (V)}
}

\commentlines{2}

\subsection{Exercise 3: Current calculation with a delta network}
\status{Modified (Ch.~7, Ex.~3: delta network with current source).}
In the circuit of Figure~\ref{fig:lab6_ex03} a 6\,mA current source feeds node A. The resistors are R1 = 6\,k$\Omega$ (A--B), R2 = 3\,k$\Omega$ (B--C), R3 = 12\,k$\Omega$ (A--C), R4 = 4\,k$\Omega$ (B to ground) and R5 = 2\,k$\Omega$ (C to ground). Use a $\Delta$--Y transformation (or nodal analysis) to calculate the three node voltages and the current in every resistor. Then simulate.

\labfig{lab6_ex03.png}{9cm}{Exercise 3: delta network with a current source}{fig:lab6_ex03}

\theory{12}

\simulation{
\item Draw the circuit of Figure~\ref{fig:lab6_ex03} in LTspice and save it as \file{Lab6/Ex03.asc}.
\item Add the \cmd{.op} analysis (\textit{Simulate $\rightarrow$ Edit Simulation Cmd $\rightarrow$ DC op pnt}), run the simulation and read the node voltages and branch currents.
\item Be careful with signs: LTspice reports the current of a voltage source as positive when it flows \emph{into} its positive terminal, and the current of a two-terminal device as positive when it flows from its first pin to its second pin.
}

\comparetable{
\cmprow{$V_A$ (V)}
\cmprow{$V_B$ (V)}
\cmprow{$V_C$ (V)}
\cmprow{$I_{R1}$ (mA)}
\cmprow{$I_{R2}$ (mA)}
\cmprow{$I_{R3}$ (mA)}
\cmprow{$I_{R4}$ (mA)}
\cmprow{$I_{R5}$ (mA)}
}

\commentlines{2}

\subsection{Exercise 4: Diode circuit 1}
\status{Unchanged (Ch.~7, Ex.~4: diode in a resistor bridge).}
In the circuit of Figure~\ref{fig:lab6_ex04} a 10\,V source feeds two branches, 2\,k$\Omega$ with 8\,k$\Omega$ and 4\,k$\Omega$ with 6\,k$\Omega$, and a diode is connected between the two middle nodes ($V_1$ at the anode, $V_2$ at the cathode). Assume the diode drops 0.7\,V when it conducts. Determine whether it conducts, then calculate $V_1$, $V_2$ and the currents through all the resistors and through the diode. Then simulate.

\labfig{lab6_ex04.png}{8cm}{Exercise 4: diode in a resistor bridge}{fig:lab6_ex04}

\theory{12}

\simulation{
\item Draw the circuit of Figure~\ref{fig:lab6_ex04} in LTspice and save it as \file{Lab6/Ex04.asc}.
\item Declare the diode model with a directive (\texttt{S} key), for example \cmd{.model D1N4148 D(Is=2.52n Rs=0.568 N=1.752)}, and type the model name in the value field of the diode.
\item Add the \cmd{.op} analysis (\textit{Simulate $\rightarrow$ Edit Simulation Cmd $\rightarrow$ DC op pnt}), run the simulation and read the node voltages and branch currents.
\item Be careful with signs: LTspice reports the current of a voltage source as positive when it flows \emph{into} its positive terminal, and the current of a two-terminal device as positive when it flows from its first pin to its second pin.
}

\comparetable{
\cmprow{$V_1$ (V)}
\cmprow{$V_2$ (V)}
\cmprow{$I_D$ (mA)}
\cmprow{$I_{R1}$ (mA)}
\cmprow{$I_{R2}$ (mA)}
\cmprow{$I_{R3}$ (mA)}
\cmprow{$I_{R4}$ (mA)}
}

\commentlines{2}

\subsection{Exercise 5: Diode circuit 2}
\status{Unchanged (Ch.~7, Ex.~5: diode with resistor network).}
In the circuit of Figure~\ref{fig:lab6_ex05} a 30\,V battery feeds R1 = 1\,k$\Omega$ in series with R2 = 1.5\,k$\Omega$; at the end of R2 the circuit splits into R3 = 4.7\,k$\Omega$ to ground and the series connection of R4 = 4.7\,k$\Omega$ and the diode D1 to ground. Decide whether the diode conducts (use $V_F=0.7$\,V), then determine the current in each resistor and in the diode and the voltage at the node between R2, R3 and R4. Then simulate.

\labfig{lab6_ex05.png}{10cm}{Exercise 5: diode with a resistor network}{fig:lab6_ex05}

\theory{12}

\simulation{
\item Draw the circuit of Figure~\ref{fig:lab6_ex05} in LTspice and save it as \file{Lab6/Ex05.asc}.
\item Declare the diode model with a directive (\texttt{S} key), for example \cmd{.model D1N4148 D(Is=2.52n Rs=0.568 N=1.752)}, and type the model name in the value field of the diode.
\item Add the \cmd{.op} analysis (\textit{Simulate $\rightarrow$ Edit Simulation Cmd $\rightarrow$ DC op pnt}), run the simulation and read the node voltages and branch currents.
\item Be careful with signs: LTspice reports the current of a voltage source as positive when it flows \emph{into} its positive terminal, and the current of a two-terminal device as positive when it flows from its first pin to its second pin.
}

\comparetable{
\cmprow{$I_{R1}$ (mA)}
\cmprow{$I_{R2}$ (mA)}
\cmprow{$I_{R3}$ (mA)}
\cmprow{$I_{R4}$ (mA)}
\cmprow{$I_{D}$ (mA)}
\cmprow{Voltage at the node (V)}
}

\commentlines{2}

\subsection{Exercise 6: Zener diode regulator}
\status{Modified (Ch.~7, Ex.~6: Zener diode voltage regulator).}
In the circuit of Figure~\ref{fig:lab6_ex06} a 16\,V source feeds, through a series resistor R = 220\,$\Omega$, a Zener diode ($V_Z=12$\,V) in parallel with a load R$_L$. Determine the load voltage and the currents $I$, $I_Z$ and $I_L$ for R$_L$ = 10\,k$\Omega$ and for R$_L$ = 1\,k$\Omega$, and the power dissipated in the Zener diode in each case. What is the minimum load resistance for which the Zener diode still regulates? Then simulate.

\labfig{lab6_ex06.png}{9cm}{Exercise 6: Zener regulator}{fig:lab6_ex06}

\theory{12}

\simulation{
\item Draw the circuit of Figure~\ref{fig:lab6_ex06} in LTspice and save it as \file{Lab6/Ex06.asc}.
\item Declare the Zener model with a directive, for example \cmd{.model DZ12 D(Is=1e-14 Bv=12 Ibv=10m Rs=2)}, and place the diode in the reverse direction (cathode to the positive node).
\item Set the load to \cmd{\{RL\}} with \cmd{.param RL=10k} and run both loads with \cmd{.step param RL list 1k 10k}.
\item Add the \cmd{.op} analysis (\textit{Simulate $\rightarrow$ Edit Simulation Cmd $\rightarrow$ DC op pnt}), run the simulation and read the node voltages and branch currents.
\item Be careful with signs: LTspice reports the current of a voltage source as positive when it flows \emph{into} its positive terminal, and the current of a two-terminal device as positive when it flows from its first pin to its second pin.
}

\comparetable{
\cmprow{$V_L$, 10\,k$\Omega$ (V)}
\cmprow{$I$, 10\,k$\Omega$ (mA)}
\cmprow{$I_L$, 10\,k$\Omega$ (mA)}
\cmprow{$I_Z$, 10\,k$\Omega$ (mA)}
\cmprow{$P_Z$, 10\,k$\Omega$ (mW)}
\cmprow{$V_L$, 1\,k$\Omega$ (V)}
\cmprow{$I$, 1\,k$\Omega$ (mA)}
\cmprow{$I_L$, 1\,k$\Omega$ (mA)}
\cmprow{$I_Z$, 1\,k$\Omega$ (mA)}
\cmprow{$P_Z$, 1\,k$\Omega$ (mW)}
\cmprow{$R_{L,min}$ ($\Omega$)}
}

\commentlines{2}

\subsection{Exercise 7: BJT driven into saturation}
\status{Modified (Ch.~7, Ex.~7: BJT driven into saturation with DC sweep).}
In the circuit of Figure~\ref{fig:lab6_ex07} an input source V1 drives the base of an NPN transistor through R1 = 10\,k$\Omega$; the collector is connected to a 10\,V supply through R2 = 470\,$\Omega$ and the emitter is grounded. Assume $\beta=100$, $V_{BE}=0.7$\,V and $V_{CE(sat)}=0.2$\,V. Determine the input voltage $V_1$ at which the transistor enters saturation and the collector current at that point. Then sweep V1 from 0 to 6\,V with a DC sweep and plot $V_C$.

\labfig{lab6_ex07.png}{9cm}{Exercise 7: BJT driven by a swept source}{fig:lab6_ex07}

\theory{10}

\simulation{
\item Draw the circuit of Figure~\ref{fig:lab6_ex07} in LTspice and save it as \file{Lab6/Ex07.asc}.
\item Declare \cmd{.model QN NPN(Is=1e-14 Bf=100 Br=1)} and type \texttt{QN} in the value field of the transistor.
\item Run \cmd{.dc V1 0 6 0.05} and plot \cmd{V(c)} and \cmd{Ic(Q1)}.
\item Find the value of $V_1$ at which the collector voltage reaches the saturation value (cursors).
}

\comparetable{
\cmprow{$V_1$ at saturation (V)}
\cmprow{$I_C$ at saturation (mA)}
\cmprow{$V_C$ for $V_1=1$\,V (V)}
\cmprow{$V_C$ for $V_1=6$\,V (V)}
}

\commentlines{2}

\subsection{Exercise 8: BJT with a Zener reference}
\status{Unchanged (Ch.~7, Ex.~8: BJT with Zener reference diode).}
In the circuit of Figure~\ref{fig:lab6_ex08} a Zener diode with $V_Z=1.7$\,V keeps the base of an NPN transistor at a fixed voltage; V$_{CC}$ = 6\,V, R$_B$ = 8.6\,k$\Omega$, R$_L$ = 400\,$\Omega$ in the collector, R$_E$ = 100\,$\Omega$ in the emitter. Assume $\beta=100$ and $V_{BE}=0.7$\,V. Determine $V_E$, $I_E$, $I_C$, $I_B$, $V_{CE}$, the Zener current and the collector-base voltage, and verify that the transistor works in the active region. Then simulate.

\labfig{lab6_ex08.png}{9cm}{Exercise 8: BJT with a Zener diode at the base}{fig:lab6_ex08}

\theory{12}

\simulation{
\item Draw the circuit of Figure~\ref{fig:lab6_ex08} in LTspice and save it as \file{Lab6/Ex08.asc}.
\item Declare \cmd{.model QN NPN(Is=1e-14 Bf=100 Br=1)} and type \texttt{QN} in the value field of the transistor.
\item Declare a low-voltage Zener model, for example \cmd{.model DZ17 D(Is=1e-14 Bv=1.7 Ibv=1m Rs=2)}.
\item Add the \cmd{.op} analysis (\textit{Simulate $\rightarrow$ Edit Simulation Cmd $\rightarrow$ DC op pnt}), run the simulation and read the node voltages and branch currents.
\item Be careful with signs: LTspice reports the current of a voltage source as positive when it flows \emph{into} its positive terminal, and the current of a two-terminal device as positive when it flows from its first pin to its second pin.
}

\comparetable{
\cmprow{$V_E$ (V)}
\cmprow{$I_E$ (mA)}
\cmprow{$I_C$ (mA)}
\cmprow{$I_B$ ($\micro$A)}
\cmprow{$I_Z$ (mA)}
\cmprow{$V_{CE}$ (V)}
\cmprow{$V_{CB}$ (V)}
}

\commentlines{2}

\subsection{Exercise 9: Diode clipper}
\status{Supplementary (diode clipper circuit for wave shaping).}
The circuit of Figure~\ref{fig:lab6_ex09} is a diode clipper: a 10\,V peak, 1\,kHz sine source feeds a 1\,k$\Omega$ resistor; the output node is connected through a diode D1 to a 3\,V reference (anode at the output) and through a diode D2 to ground (cathode at the output). With $V_F=0.7$\,V determine the maximum and the minimum output voltage and sketch the output. Then run a transient simulation of 4\,ms.

\labfig{lab6_ex09.png}{9cm}{Exercise 9: diode clipper}{fig:lab6_ex09}

\theory{8}

\simulation{
\item Draw the circuit of Figure~\ref{fig:lab6_ex09} in LTspice and save it as \file{Lab6/Ex09.asc}.
\item Declare the diode model with a directive (\texttt{S} key), for example \cmd{.model D1N4148 D(Is=2.52n Rs=0.568 N=1.752)}, and type the model name in the value field of the diode.
\item Set the source to SINE (amplitude 10\,V, 1\,kHz) and run \cmd{.tran 0 4m 0 5u}; plot the input and the output.
\item Read the maximum and minimum of the output with \cmd{.meas TRAN Vmax MAX V(out)} and \cmd{.meas TRAN Vmin MIN V(out)} or with the cursors.
}

\comparetable{
\cmprow{Maximum output voltage (V)}
\cmprow{Minimum output voltage (V)}
}

\commentlines{2}

\subsection{Exercise 10: Series regulator with a Zener diode and a BJT}
\status{Supplementary (series voltage regulator with Zener reference and BJT pass transistor).}
In the series regulator of Figure~\ref{fig:lab6_ex10} a 12\,V source feeds a Zener diode ($V_Z=5.1$\,V) through R1 = 1\,k$\Omega$; the Zener voltage is applied to the base of an NPN transistor (collector connected to 12\,V) whose emitter is the output, loaded by R$_L$. Assume $\beta=100$ and $V_{BE}=0.7$\,V. Determine the output voltage, the Zener current, the load current and the base current for R$_L$ = 100\,$\Omega$. Explain how the regulator keeps the output constant when the load changes from 50\,$\Omega$ to 1\,k$\Omega$ and why the transistor makes it better than a simple Zener regulator. Then simulate.

\labfig{lab6_ex10.png}{9cm}{Exercise 10: series regulator}{fig:lab6_ex10}

\theory{12}

\simulation{
\item Draw the circuit of Figure~\ref{fig:lab6_ex10} in LTspice and save it as \file{Lab6/Ex10.asc}.
\item Declare \cmd{.model QN NPN(Is=1e-14 Bf=100 Br=1)} and type \texttt{QN} in the value field of the transistor.
\item Declare \cmd{.model DZ51 D(Is=1e-14 Bv=5.1 Ibv=10m Rs=2)}.
\item Set the load to \cmd{\{RL\}}, define \cmd{.param RL=100} and run \cmd{.step param RL list 50 100 1k}.
\item Add the \cmd{.op} analysis (\textit{Simulate $\rightarrow$ Edit Simulation Cmd $\rightarrow$ DC op pnt}), run the simulation and read the node voltages and branch currents.
\item Be careful with signs: LTspice reports the current of a voltage source as positive when it flows \emph{into} its positive terminal, and the current of a two-terminal device as positive when it flows from its first pin to its second pin.
}

\comparetable{
\cmprow{$V_{out}$, $R_L=50\,\Omega$ (V)}
\cmprow{$V_{out}$, $R_L=100\,\Omega$ (V)}
\cmprow{$V_{out}$, $R_L=1\,$k$\Omega$ (V)}
\cmprow{$I_Z$, $R_L=100\,\Omega$ (mA)}
\cmprow{$I_L$, $R_L=100\,\Omega$ (mA)}
\cmprow{$I_B$, $R_L=100\,\Omega$ ($\micro$A)}
}

\commentlines{2}

\newpage
\section{Advanced exercises}

\subsection{Advanced Exercise 1: Monte Carlo tolerance analysis}
\status{Advanced / Supplementary (Monte Carlo tolerance analysis of series regulator using \cmd{mc()} and \cmd{.step}).}
The resistors of the regulator of Exercise~10 (Figure~\ref{fig:lab6_adv01}) have a 5\,\% tolerance. Study the spread of the output voltage over 20 random draws of the resistor values.

\labfig{lab6_adv01.png}{9cm}{Advanced Exercise 1: series regulator with tolerance}{fig:lab6_adv01}

\subsubsection{Theory calculation}\vspace{-3pt}\notesblock Which resistor affects the output voltage more strongly, R1 or R$_L$? Estimate the worst-case spread of $V_{out}$.\dotlines{4}

\simulation{
\item Draw the circuit of Figure~\ref{fig:lab6_adv01} in LTspice and save it as \file{Lab6/Ex11.asc}.
\item Declare the models \cmd{.model QN NPN(Is=1e-14 Bf=100 Br=1)} and \cmd{.model DZ51 D(Is=1e-14 Bv=5.1 Ibv=10m Rs=2)}.
\item Give the resistors random values with the function \cmd{mc}: R1 = \cmd{\{mc(1k,0.05)\}} and R$_L$ = \cmd{\{mc(100,0.05)\}}.
\item Add \cmd{.step param run 1 20 1} to repeat the simulation 20 times and \cmd{.op}.
\item Read the output voltage of the 20 runs in the SPICE Error Log (for example with \cmd{.meas OP Vout FIND V(out)}) and fill in the table.
}
\begin{center}\renewcommand{\arraystretch}{1.8}\begin{tabular}{|c|c|}\hline \textbf{Quantity} & \textbf{Value (V)}\\ \hline Minimum $V_{out}$ & \\ \hline Maximum $V_{out}$ & \\ \hline Average $V_{out}$ & \\ \hline Spread (max $-$ min) & \\ \hline \end{tabular}\end{center}
\commentlines{2}
\newpage

\subsection{Advanced Exercise 2: Thevenin equivalent with the \texttt{.tf} command}
\status{Advanced / Supplementary (Thevenin equivalent circuit extraction using the \cmd{.tf} transfer function command).}
Remove R5 from the bridge network of Exercise~1 and consider the terminals where R5 was connected (Figure~\ref{fig:lab6_adv02}). Find the Thevenin equivalent seen from these terminals and use it to calculate the current in R5 = 40\,$\Omega$.

\labfig{lab6_adv02.png}{9cm}{Advanced Exercise 2: network seen by a load resistor}{fig:lab6_adv02}

\subsubsection{Theory calculation}\vspace{-3pt}\notesblock Determine $V_{th}$ (open-circuit voltage) and $R_{th}$ (with the source replaced by a short circuit), then the current in a 40\,$\Omega$ load.\dotlines{6}

\simulation{
\item Draw the circuit of Figure~\ref{fig:lab6_adv02} in LTspice and save it as \file{Lab6/Ex12.asc}.
\item Run the directive \cmd{.tf V(r) V1} (output node r, input source V1). The SPICE Error Log lists the gain, the input impedance and the output impedance at the node.
\item $V_{th}=\text{gain}\times V_1$ and $R_{th}$ is the output impedance. Check the result by connecting the 40\,$\Omega$ resistor and running \cmd{.op}.
}
\begin{center}\renewcommand{\arraystretch}{1.8}\begin{tabular}{|c|c|c|c|}\hline \textbf{Quantity} & \textbf{Theory} & \textbf{LTspice} & \textbf{Error (\%)}\\ \hline $V_{th}$ (V) & & & \\ \hline $R_{th}$ ($\Omega$) & & & \\ \hline Load current (mA) & & & \\ \hline \end{tabular}\end{center}
\commentlines{2}
\newpage
